\chapter*{Course Structure}
\addcontentsline{toc}{chapter}{Course Structure}

The course develops the probability, statistics, linear algebra, calculus,
optimization, and numerical computing foundations used throughout machine
learning. The detailed programme describes thirteen substantive lecture units.
Its original numbering uses ``Slot 7'' twice; the sequence below normalizes the
labels to Lectures 1--13 without removing any listed material.

\section*{Block I: Probability and Statistics}

\begin{tabularx}{\textwidth}{@{}L{1.7cm}Y@{}}
\toprule
\textbf{Lecture} & \textbf{Topic} \\
\midrule
1 & Probability models, conditioning, and random variables \\
2 & Distributions, expectation, and variance \\
3 & Joint distributions, dependence, and random vectors \\
4 & Samples, concentration, and limit theorems \\
5 & Statistical estimation: method of moments, maximum likelihood, and MAP \\
6 & Sampling distributions, confidence intervals, and bootstrap \\
7 & Entropy, cross-entropy, KL divergence, Fisher information, and the Cram\'er--Rao bound \\
\bottomrule
\end{tabularx}

\section*{Block II: Linear Algebra, Calculus, and Optimization}

\begin{tabularx}{\textwidth}{@{}L{1.7cm}Y@{}}
\toprule
\textbf{Lecture} & \textbf{Topic} \\
\midrule
8 & Geometry of data, projections, and least squares \\
9 & Eigenvalues, positive-definite matrices, singular value decomposition, and PCA \\
10 & Multivariable calculus, matrix calculus, and automatic differentiation \\
11 & Optimization for machine learning \\
\bottomrule
\end{tabularx}

\section*{Block III: Numerical Computing and Autodiff}

\begin{tabularx}{\textwidth}{@{}L{1.7cm}Y@{}}
\toprule
\textbf{Lecture} & \textbf{Topic} \\
\midrule
12 & NumPy, visualization, and numerical linear algebra \\
13 & PyTorch, autograd, and the integrated capstone \\
\bottomrule
\end{tabularx}

The notes are organized as one cumulative conspect: each future lecture will
extend the same PDF.
